\documentclass[10pt,a4paper]{article}
\usepackage{amsmath}
\usepackage{graphicx}
\usepackage{booktabs}
\usepackage{float}
\usepackage[margin=0.75in]{geometry}

\title{\textbf{Project Phase 3: RoboMaster Odometry}\\[0.2em]
\normalsize Advanced Robotics (SCE-672) -- King Fahd University of Petroleum \& Minerals}
\author{Raian E Rabbani (202617880) \\
\small Instructor: Prof. Aamer Iqbal Bhatti}
\date{November 2026}

\begin{document}
\maketitle

\section{Approach and Sensor Measurements}
The system was implemented on Ubuntu using ROS~2 with the DJI RoboMaster EP Core. Rather than sending high-level $(v, \omega)$ chassis velocity commands, individual wheel speeds were commanded directly in RPM using Mecanum inverse kinematics (\texttt{drive\_wheels}). A two-state Kalman filter $\mathbf{x}_k = [d_k, v_k]^T$ tracks the bumper-to-tag distance and forward velocity:
\begin{itemize}
    \item \textbf{Tag Detection:} Frames from the 720p camera detect a fixed 7.8~cm AprilTag (\texttt{36h11}). Distance is obtained from calibrated pinhole geometry ($f_x = 688.6$~px) minus the camera-to-bumper offset ($d_{\text{off}} = 0.1876$~m): $z_k = \frac{f_x W_{\text{tag}}}{w_{\text{px}}} - d_{\text{off}}$.
    \item \textbf{Straight-Line Motion:} A PD controller on the visual bearing angle keeps the robot aligned with the tag while forward speed ramps down from $0.28$~m/s to $0.105$~m/s before stopping.
    \item \textbf{Prediction Sensor (50~Hz):} Chassis wheel encoder odometry supplies incremental displacement $\Delta x_{\text{odo}}$ to predict $\hat{d}_{k|k-1} = \hat{d}_{k-1|k-1} - \Delta x_{\text{odo}}$ between camera updates.
    \item \textbf{Correction Sensor:} Each camera tag distance $z_k$ corrects the predicted state and covariance to compensate for wheel slip.
\end{itemize}

\section{Experimental Setup}
All 10 repeat trials were run from a starting distance of 1.50~m with a target stopping distance of 0.25~m. Both the estimated and measured distances were referenced to the robot's front bumper.

\section{Results and Error Calculations}
Estimation error is calculated as $e = \hat{d} - d_m$ (estimated distance minus measured distance). Table~\ref{tab:results} summarizes all 10 trials alongside pure wheel odometry for comparison.

\begin{table}[H]
\centering
\caption{Final Distance Estimates and Errors Across 10 Repeat Trials ($d_{\text{stop}} = 0.25$~m)}
\label{tab:results}
\small
\begin{tabular}{ccccccc}
\toprule
\textbf{Trial} & \textbf{Actual $d_m$ (m)} & \textbf{KF Est $\hat{d}$ (m)} & \textbf{KF Error $e$ (m)} & \textbf{$|e|$ (m)} & \textbf{Pure Odo (m)} & \textbf{Odo Error (m)} \\
\midrule
1 & 0.2596 & 0.2656 & +0.0060 & 0.0060 & 0.2768 & +0.0172 \\
2 & 0.2543 & 0.2623 & +0.0080 & 0.0080 & 0.2848 & +0.0305 \\
3 & 0.2551 & 0.2639 & +0.0088 & 0.0088 & 0.2979 & +0.0428 \\
4 & 0.2567 & 0.2653 & +0.0086 & 0.0086 & 0.3776 & +0.1209 \\
5 & 0.2520 & 0.2604 & +0.0084 & 0.0084 & 0.2700 & +0.0180 \\
6 & 0.2503 & 0.2583 & +0.0080 & 0.0080 & 0.3036 & +0.0533 \\
7 & 0.2519 & 0.2598 & +0.0079 & 0.0079 & 0.2943 & +0.0424 \\
8 & 0.2639 & 0.2639 & +0.0000 & 0.0000 & 0.3409 & +0.0770 \\
9 & 0.2607 & 0.2607 & -0.0000 & 0.0000 & 0.3980 & +0.1373 \\
10 & 0.2689 & 0.2689 & +0.0000 & 0.0000 & 0.3087 & +0.0398 \\
\midrule
\textbf{MAE} & -- & -- & -- & \textbf{0.0056} & -- & 0.0579 \\
\textbf{Range} & -- & -- & [-0.0000, +0.0088] & \textbf{0.0088} & -- & 0.1201 \\
\bottomrule
\end{tabular}
\end{table}

\begin{figure}[H]
\centering
\includegraphics[width=0.84\textwidth]{/home/raian/Documents/phd/adv robotics/project/results/plots/trial_error_comparison.pdf}
\caption{Scatterplot of final distance estimation errors across 10 trials.}
\end{figure}

\section{Discussion}
Across the 10 trials, the Kalman filter produced consistent stops near 0.25~m with a Mean Absolute Error of 0.56~cm and an error range of 0.88~cm. By comparison, pure wheel odometry had an MAE of 5.79~cm and consistently under-counted the traveled distance.

The primary sources of error were Mecanum roller slip during acceleration and stopping, minor floor unevenness, motor braking inertia, and pixel-level corner quantization on the 7.8~cm tag. Odometry prediction provided smooth 50~Hz position tracking between camera frames, while visual correction removed the accumulated slip error so the robot stopped reliably at the target distance.

\end{document}
